\documentclass[12pt]{article}
\usepackage[a4paper,margin=1in]{geometry}
\usepackage[utf8]{inputenc}
\usepackage[T1]{fontenc}
\usepackage{amsmath,amssymb}
\usepackage{textcomp}
\usepackage{parskip}
\begin{document}
\begin{center}
{\Large \textbf{2015 STEP I}}
\end{center}

\subsection*{Page 1}
\begin{flushleft}
\bigskip\textbf{Section A: Pure Mathematics}\par
x 2\par
\bigskip\textbf{Question 1.} (i) Sketch the curve y = e (2x 5x + 2) .\par
-\par
x 2\par
\[Hence determine how many real values of x satisfy the equation e (2x 5x + 2) = k\]
-\par
in the different cases that arise according to the value of k.\par
n x\par
You may assume that x e 0 as x for any integer n.\par
\to \to -\infty\par
\[x^2 4 2\]
\textbf{(ii)} Sketch the curve y = e (2x 5x + 2) .\par
-\par
\[\sqrt\{3\} + 1\]
\textasciicircum{}\circ \textasciicircum{}\circ\par
\bigskip\textbf{Question 2.} (i) Show that cos 15 = and find a similar expression for sin 15 .\par
\[2\sqrt\{2\}\]
\textbf{(ii)} Show that cos α is a root of the equation\par
\[4x 3x cos 3α = 0 ,\]
- -\par
and find the other two roots in terms of cos α and sin α.\par
\textbf{(iii)} Use parts (i) and (ii) to solve the equation y 3 3y \sqrt2 = 0 , giving your answers in\par
- -\par
surd form.\par
\bigskip\textbf{Question 3.} A prison consists of a square courtyard of side b bounded by a perimeter wall and a square\par
building of side a placed centrally within the courtyard. The sides of the building are parallel\par
to the perimeter walls.\par
Guards can stand either at the middle of a perimeter wall or in a corner of the courtyard. If\par
the guards wish to see as great a length of the perimeter wall as possible, determine which\par
of these positions is preferable. You should consider separately the cases b < 3a and b > 3a .\par
\end{flushleft}
\subsection*{Page 2}
\begin{flushleft}
\bigskip\textbf{Question 4.} The midpoint of a rod of length 2b slides on the curve y = 1 x 2 , x \leq 0, in such a way that\par
the rod is always tangent, at its midpoint, to the curve. Show that the curve traced out by\par
one end of the rod can be written in the form\par
\[x = 2 tan θ b cos θ\]
-\par
\[y = tan θ b sin θ\]
-\par
\[for some suitably chosen angle θ which satisfies 0 \geq θ < 1 \pi .\]
When one end of the rod is at a point A on the y-axis, the midpoint is at point P and θ = α.\par
Let R be the region bounded by the following:\par
1 2\par
\[the curve y = x between the origin and P ;\]
the y-axis between A and the origin;\par
the half-rod AP .\par
2 3\par
Show that the area of R is tan α.\par
\bigskip\textbf{Question 5.} (i) The function f is defined, for x > 0, by\par
x-1\par
\[f(x) = (t 1) dt .\]
\[\leq 1 -\]
\[By evaluating the integral, sketch the curve y = f(x).\]
\textbf{(ii)} The function g is defined, for < x < , by\par
-\infty \infty\par
\[g(x) = dt .\]
\[\leq -1 \sqrt\{1\} 2xt + x^2\]
-\par
\[By evaluating the integral, sketch the curve y = g(x).\]
\[\geq \geq\]
\bigskip\textbf{Question 6.} The vertices of a plane quadrilateral are labelled A, B, A and B , in clockwise order. A\par
\[\geq \geq\]
point O lies in the same plane and within the quadrilateral. The angles AOB and A OB\par
\[\geq \geq\]
\[are right angles, and OA = OB and OA = OB .\]
\[\geq \geq \geq \geq\]
Use position vectors relative to O to show that the midpoints of AB, BA , A B and B A\par
are the vertices of a square.\par
\[\geq\]
Given that the lengths of OA and OA are fixed (and the conditions of the first paragraph\par
\[\geq\]
still hold), find the value of angle BOA for which the area of the square is greatest.\par
\end{flushleft}
\subsection*{Page 3}
\begin{flushleft}
7 Let\par
2 3\par
\[f(x) = 3ax 6x\]
-\par
and, for each real number a, let M(a) be the greatest value of f(x) in the interval 1 \geq x \geq 1.\par
-\par
Determine M(a) for a \leq 0. [The formula for M(a) is different in different ranges of a; you\par
will need to identify three ranges.]\par
\bigskip\textbf{Question 8.} Show that:\par
\textbf{(i)} 1 + 2 + 3 + + n = n(n + 1);\par
\cdots\par
\textbf{(ii)} if N is a positive integer, m is a non-negative integer and k is a positive odd integer,\par
k k\par
then (N m) + m is divisible by N .\par
-\par
k k k k\par
\[Let S = 1 + 2 + 3 + + n , where k is a positive odd integer. Show that if n is odd then\]
\cdots 1\par
S is divisible by n and that if n is even then S is divisible by n.\par
Show further that S is divisible by 1 + 2 + 3 + + n.\par
\cdots\par
\end{flushleft}
\subsection*{Page 4}
\begin{flushleft}
\bigskip\textbf{Section B: Mechanics}\par
\bigskip\textbf{Question 9.} A short-barrelled machine gun stands on horizontal ground. The gun fires bullets, from\par
\[\pi\]
\[ground level, at speed u continuously from t = 0 to t = , where λ is a positive constant,\]
6λ\par
but does not fire outside this time period. During this time period, the angle of elevation α\par
1 1\par
\[of the barrel decreases from \pi to \pi and is given at time t by\]
3 6\par
\[α = \pi λt .\]
-\par
g\par
Let k = . Show that, in the case 1 \geq k \geq 1 \sqrt3, the last bullet to hit the ground does so\par
2λu 2 2\par
at a distance\par
\[2ku 2\sqrt\{1\} k^2\]
-\par
g\par
from the gun.\par
\[What is the corresponding result if k < ?\]
\bigskip\textbf{Question 10.} A bus has the shape of a cuboid of length a and height h. It is travelling northwards on a\par
journey of fixed distance at constant speed u (chosen by the driver). The maximum speed of\par
the bus is w. Rain is falling from the southerly direction at speed v in straight lines inclined\par
\[to the horizontal at angle θ, where 0 < θ < \pi.\]
By considering first the case u = 0, show that for u > 0 the total amount of rain that hits\par
the roof and the back or front of the bus in unit time is proportional to\par
h v cos θ u + av sin θ .\par
-\par
\[\geq \geq\]
\[\geq \geq\]
Show that, in order to encounter as little rain as possible on the journey, the driver should\par
\[choose u = w if either w < v cos θ or a sin θ > h cos θ. How should the speed be chosen if\]
\[w > v cos θ and a sin θ < h cos θ? Comment on the case a sin θ = h cos θ.\]
How should the driver choose u on the return journey?\par
\end{flushleft}
\subsection*{Page 5}
\begin{flushleft}
\bigskip\textbf{Question 11.} Two long circular cylinders of equal radius lie in equilibrium on an inclined plane, in contact\par
with one another and with their axes horizontal. The weights of the upper and lower\par
cylinders are W1 and W2, respectively, where W1 > W2 . The coefficients of friction between\par
the inclined plane and the upper and lower cylinders are μ1 and μ2, respectively, and the\par
coefficient of friction between the two cylinders is μ. The angle of inclination of the plane\par
is α (which is positive).\par
\textbf{(i)} Let F be the magnitude of the frictional force between the two cylinders, and let F1\par
\[and F^2 be the magnitudes of the frictional forces between the upper cylinder and the\]
\[plane, and the lower cylinder and the plane, respectively. Show that F = F^1 = F^2 .\]
\textbf{(ii)} Show that\par
\[W^1 + W^2\]
\[μ \leq ,\]
\[W^1 W^2\]
-\par
and that\par
\[2μ1W1\]
\[tan α \geq .\]
\[(1 + μ1)(W^1 + W^2)\]
\end{flushleft}
\subsection*{Page 6}
\begin{flushleft}
\bigskip\textbf{Section C: Probability and Statistics}\par
\bigskip\textbf{Question 12.} The number X of casualties arriving at a hospital each day follows a Poisson distribution\par
with mean 8; that is,\par
-8 n\par
e 8\par
\[P(X = n) = , n = 0, 1, 2, . . . .\]
n!\par
Casualties require surgery with probability . The number of casualties arriving on any\par
given day is independent of the number arriving on any other day and the casualties require\par
surgery independently of one another.\par
\textbf{(i)} What is the probability that, on a day when exactly n casualties arrive, exactly r of\par
them require surgery?\par
\textbf{(ii)} Prove (algebraically) that the number requiring surgery each day also follows a Poisson\par
distribution, and state its mean.\par
\textbf{(iii)} Given that in a particular randomly chosen week a total of 12 casualties require surgery\par
on Monday and Tuesday, what is the probability that 8 casualties require surgery on\par
Monday? You should give your answer as a fraction in its lowest terms.\par
\bigskip\textbf{Question 13.} A fair die with faces numbered 1, . . . , 6 is thrown repeatedly. The events A, B, C, D and E\par
are defined as follows.\par
A: the first 6 arises on the nth throw.\par
B: at least one 5 arises before the first 6.\par
C: at least one 4 arises before the first 6.\par
D: exactly one 5 arises before the first 6.\par
E: exactly one 4 arises before the first 6.\par
Evaluate the following probabilities:\par
\textbf{(i)} P(A) (ii) P(B) (iii) P(B C) (iv) P(D) (v) P(D E) .\par
\cap \cup\par
For some parts of this question, you may want to make use of the binomial expansion in the\par
form:\par
n(n + 1) (n + r 1)!\par
-n 2 r\par
\[(1 x) = 1 + nx + x + + - x + .\]
- 2 \cdots r!(n 1)! \cdots\par
-\par
\end{flushleft}
\end{document}
